\Needspace{8\baselineskip}
\Article{258}
Alle im vorigen Artikel gefundenen Formen gehören offenbar zu ambigen Klassen, und umgekehrt muss jede eigentlich primitive ambige Klasse der Determinante \(D\) mindestens eine jener Formen enthalten. Denn in einer solchen Klasse gibt es gewiss ambige Formen, und jeder eigentlich primitiven ambigen Form \((a,b,c)\) der Determinante \(D\) ist eine der Formen des vorigen Artikels äquivalent, nämlich entweder
\[
(a,0,-\tfrac{D}{a})\text{ oder }(a,\tfrac12 a,\tfrac14 a-\tfrac{D}{a})
\]
je nachdem, ob \(b\) entweder \(\equiv0\) oder \(\equiv\frac12 a\ (\mathrm{mod.}\,a)\) ist. Die Aufgabe ist also darauf zurückgeführt, zu untersuchen, wie viele verschiedene Klassen jene Formen bilden.

Ist die Form \((a,0,c)\) unter den Formen des vorigen Artikels, so kommt auch die Form \((c,0,a)\) darunter vor und ist stets von ihr verschieden, mit der einzigen Ausnahme \(a=c=\pm1\), also \(D=-1\), die wir vorerst beiseitelassen. Da diese Formen offenbar derselben Klasse angehören, genügt es, eine zu behalten; wir entfernen diejenige, deren erster Koeffizient größer als der dritte ist. Auch den Fall \(a=-c=\pm1\) oder \(D=1\) lassen wir beiseite. So können wir alle Formen
\SourceLeaf{292}
$(A,0,C)$ auf die Hälfte ihrer Anzahl reduzieren, indem wir von je zweien stets eine behalten; in allen übrigen gilt $A<\sqrt{\pm D}$.

Kommt ebenso unter den Formen des vorigen Artikels die Form $(2b,b,c)$ vor, so findet sich darunter auch
\[
(4c-2b,\ 2c-b,\ c)=\left(-\frac{2D}{b},\ -\frac{D}{b},\ c\right)
\]
die ihr eigentlich äquivalent und von ihr verschieden ist, mit der einzigen Ausnahme $c=b=\pm1$ oder $D=-1$, die wir beiseitelassen. Von diesen beiden Formen genügt es, diejenige zu behalten, deren erster Koeffizient kleiner als der erste Koeffizient der anderen ist (dem Betrag nach gleich und im Vorzeichen verschieden können sie in diesem Fall nicht sein). Somit lassen sich auch alle Formen $(2B,B,C)$ auf die Hälfte ihrer Anzahl reduzieren, indem man von je zweien stets eine entfernt; in den übrigen gilt $B<\frac{D}{B}$ oder $B<\sqrt{\pm D}$. So bleibt von allen Formen des vorigen Artikels nur die Hälfte übrig. Ihre Menge bezeichnen wir mit $W$, und es bleibt nur zu zeigen, wie viele verschiedene Klassen aus diesen Formen entstehen. Ist $D$ negativ, so gibt es offenbar in $W$ ebenso viele positive wie negative Formen.

I. Ist $D$ negativ, so gehören die einzelnen Formen in $W$ zu verschiedenen Klassen. Denn alle Formen $(A,0,C)$ sind reduziert; ebenso sind alle Formen $(2B,B,C)$ reduziert, außer denen, in denen $C<2B$ gilt. In einer solchen Form ist $2C<2B+C$; daher gilt (weil $B<\frac{D}{B}$, \textit{das heißt}\ $B<2C-B$, also $2B<2C$ oder $B<C$) $2C-2B<C$ und $C-B<\frac12 C$, und somit ist $(C,C-B,C)$, die jener offenbar äquivalent ist, eine reduzierte Form. So erhält man ebenso viele reduzierte Formen, wie es Formen in $W$ gibt. Da darunter offenbar weder identische noch entgegengesetzte Formen vorkommen können (mit der einzigen Ausnahme $C-B=0$, bei der $B=C=\pm1$, also $D=-1$ gilt und die wir bereits beiseitegelassen haben), gehören alle zu verschiedenen Klassen. Daraus folgt: Die Anzahl aller eigentlich primitiven ambigen Klassen der Determinante $D$ ist gleich der Anzahl der Formen in $W$ oder der Hälfte der Anzahl der Formen des vorigen Artikels. Im Ausnahmefall $D=-1$ ergibt sich durch Ausgleich dasselbe: Es gibt zwei Klassen; zur einen gehören die Formen $(1,0,1)$, $(2,1,1)$, zur anderen $(-1,0,-1)$, $(-2,-1,-1)$. Allgemein ist also bei negativer Determinante die Anzahl aller eigentlich primitiven ambigen Klassen gleich der Anzahl aller
\SourceLeaf{293}
zuweisbaren Charaktere primitiver Formen dieser Determinante; die Anzahl der positiven eigentlich primitiven ambigen Klassen ist dagegen halb so groß.

II. Ist $D$ eine positive Quadratzahl $=hh$, so lässt sich ohne große Mühe beweisen, dass die einzelnen Formen in $W$ zu verschiedenen Klassen gehören. In diesem Fall können wir aber auf folgende Weise noch kürzer zur Lösung gelangen. Nach Artikel 210 enthält jede eigentlich primitive ambige Klasse der Determinante $hh$ und keine andere Klasse genau eine reduzierte Form $(a,h,0)$, in der $a$ ein Wert des Ausdrucks $\sqrt{1}\ (\mathrm{mod.}\ 2h)$ ist, der zwischen $0$ und $2h-1$ einschließlich der Grenzen liegt. Es gibt also so viele eigentlich primitive ambige Klassen der Determinante $hh$, wie dieser Ausdruck Werte hat. Aus Artikel 105 folgt ohne Mühe, dass die Anzahl dieser Werte $2^n$ oder $2^{n+1}$ oder $2^{n+2}$ ist, je nachdem, ob $h$ ungerade, das Doppelte einer ungeraden Zahl oder durch vier teilbar ist, also je nachdem, ob $D\equiv1$ oder $\equiv4$ oder $\equiv0\ (\mathrm{mod.}\ 8)$ gilt; dabei bezeichnet $n$ die Anzahl der verschiedenen ungeraden Primteiler von $h$ oder von $D$. Somit ist die Anzahl der eigentlich primitiven ambigen Klassen stets die Hälfte der Gesamtzahl aller im vorigen Artikel gefundenen Formen oder gleich der Anzahl der Formen in $W$ oder aller möglichen Charaktere.

III. Ist $D$ eine positive Nichtquadratzahl, so leiten wir aus jeder Form $(A,B,C)$ in $W$ eine andere Form $(A,B',C')$ ab, indem wir $B'\equiv B\ (\mathrm{mod.}\ A)$ zwischen den Grenzen $\sqrt{D}$ und $\sqrt{D}\mp A$ nehmen (mit dem oberen oder unteren Vorzeichen, je nachdem, ob $A$. positiv oder negativ ist) und $C'=\frac{B'B'-D}{A}$ setzen; die Menge dieser Formen bezeichnen wir mit $W'$. Offenbar sind sie eigentlich primitive ambige Formen der Determinante $D$ und sämtlich voneinander verschieden; außerdem sind sie alle reduziert. Gilt nämlich $A<\sqrt{D}$, so ist $B'$ offenbar $<\sqrt{D}$ und positiv; ferner gilt $B'>\sqrt{D}\mp A$, also $A>\sqrt{D}-B'$, und somit liegt $A$, positiv genommen, gewiss zwischen $\sqrt{D}+B'$ und $\sqrt{D}-B'$. Gilt dagegen $A>\sqrt{D}$, so kann nicht $B=0$ sein (solche Formen haben wir ja entfernt), sondern es muss $B=\frac12 A$ gelten; daher ist $B'$ dem Betrag nach gleich $\frac12 A$ und dem Vorzeichen nach positiv (denn wegen $A<2\sqrt{D}$ liegt $\pm\frac12 A$ zwischen den für $B'$ festgelegten Grenzen und ist $B$ modulo $A$ kongruent, also $B'=\pm\frac12 A$). Somit gilt $B'<\sqrt{D}$, woraus $2B'<\sqrt{D}+B'$ oder $A<\sqrt{D}+B'$ folgt; daher liegt $\pm A$ notwendig zwischen den Grenzen $\sqrt{D}+B'$ und $\sqrt{D}-B'$. Schließlich enthält $W'$ alle reduzierten eigentlich primitiven ambigen Formen der Determinante $D$. Ist nämlich $(a,b,c)$ eine solche Form, so gilt entweder $b\equiv0$ oder $b\equiv\frac12 a\ (\mathrm{mod.}\ a)$. Im ersten Fall kann offenbar nicht $b<a$ und somit auch nicht $a>\sqrt{D}$ gelten, weshalb die Form $\left(a,0,-\frac{D}{a}\right)$ gewiss
\SourceLeaf{294}
in $W$ und die entsprechende Form $(a,b,c)$ in $W'$ enthalten ist; im zweiten Fall gilt gewiss $a<2\sqrt{D}$, also liegt $\left(a,\allowbreak \frac12 a,\allowbreak \frac14 a-\frac{D}{a}\right)$ in $W$ und die entsprechende Form $(a,b,c)$ in $W'$. Daraus folgt, dass die Anzahl der Formen in $W$ gleich der Anzahl aller reduzierten eigentlich primitiven ambigen Formen der Determinante $D$ ist. Da aber jede ambige Klasse \textit{zwei} reduzierte ambige Formen enthält (Artikel 187, 194), ist die Anzahl aller eigentlich primitiven ambigen Klassen der Determinante $D$ halb so groß wie die Anzahl der Formen in $W$ oder wie die Anzahl aller zuweisbaren Charaktere.

\Needspace{8\baselineskip}
\Article{259}
Die Anzahl der uneigentlich primitiven ambigen Klassen einer gegebenen Determinante $D$ ist stets gleich der Anzahl der eigentlich primitiven ambigen Klassen derselben Determinante. Es sei $K$ die Hauptklasse, und $K'$, $K''$ usw. seien die übrigen eigentlich primitiven ambigen Klassen dieser Determinante; $L$ sei irgendeine uneigentlich primitive ambige Klasse derselben Determinante, \textit{zum Beispiel}\ diejenige, in der die Form $(2,1,\frac12-\frac12 D)$ liegt. Durch Komposition der Klasse $L$ mit $K$ entsteht also die Klasse $L$ selbst. Aus der Komposition der Klasse $L$ mit $K'$, $K''$ usw. mögen jeweils die Klassen $L'$, $L''$ usw. entstehen, die offenbar alle zur selben Determinante $D$ gehören und uneigentlich primitiv und ambig sind. Der Satz ist also bewiesen, sobald gezeigt ist, dass alle Klassen $L$, $L'$, $L''$ usw. verschieden sind und es außer ihnen keine uneigentlich primitiven ambigen Klassen der Determinante $D$ gibt. Dazu unterscheiden wir die folgenden Fälle.

I. Ist die Anzahl der uneigentlich primitiven Klassen gleich der Anzahl der eigentlich primitiven, so entsteht jede der ersten durch Komposition der Klasse $L$ mit einer bestimmten eigentlich primitiven Klasse; daher sind notwendig alle $L$, $L'$, $L''$ usw. verschieden. Bezeichnet nun $\mathfrak{L}$ irgendeine uneigentlich primitive ambige Klasse der Determinante $D$, so gibt es eine eigentlich primitive Klasse $\mathfrak{K}$ mit $\mathfrak{K}+L=\mathfrak{L}$. Ist der Klasse $\mathfrak{K}$ die Klasse $\mathfrak{K}'$ entgegengesetzt, so gilt auch (weil die Klassen $L$, $\mathfrak{L}$ sich selbst entgegengesetzt sind) $\mathfrak{K}'+L=\mathfrak{L}$; daher ist notwendig $\mathfrak{K}$ mit $\mathfrak{K}'$ identisch und somit ambig. Also findet sich $\mathfrak{K}$ unter den Klassen $K$, $K'$, $K''$ usw. und $\mathfrak{L}$ unter $L$, $L'$, $L''$ usw.

II. Ist die Anzahl der uneigentlich primitiven Klassen ein Drittel der Anzahl der eigentlich primitiven, so sei $H$ die Klasse, in der die Form $\left(4,\allowbreak 1,\allowbreak \frac{1-D}{4}\right)$ liegt, und $H'$ diejenige, in der $\left(4,\allowbreak 3,\allowbreak \frac{9-D}{4}\right)$ liegt; dann sind $H$, $H'$ eigentlich primitiv und sowohl
\SourceLeaf{295}
voneinander als auch von der Hauptklasse $K$ verschieden, und es gilt $H+H'=K$; $2H=H'$, $2H'=H$. Ist $\mathfrak{L}$ irgendeine uneigentlich primitive Klasse der Determinante $D$, die durch Komposition der Klasse $L$ mit der eigentlich primitiven Klasse $\mathfrak{K}$ entsteht, so gilt auch $\mathfrak{L}=L+\mathfrak{K}+H$ und $\mathfrak{L}=L+\mathfrak{K}+H'$. Außer den drei (eigentlich primitiven und verschiedenen) Klassen $\mathfrak{K}$, $\mathfrak{K}+H$, $\mathfrak{K}+H'$ gibt es keine anderen, die mit $L$ zusammengesetzt $\mathfrak{L}$ ergeben. Ist also $\mathfrak{L}$ ambig und $\mathfrak{K}'$ der Klasse $\mathfrak{K}$ entgegengesetzt, so gilt auch $L+\mathfrak{K}'=\mathfrak{L}$; daher muss $\mathfrak{K}'$ mit einer jener drei Klassen identisch sein. Gilt $\mathfrak{K}'=\mathfrak{K}$, so ist $\mathfrak{K}$ ambig; gilt $\mathfrak{K}'=\mathfrak{K}+H$, so gilt $K=\mathfrak{K}+\mathfrak{K}'=2\mathfrak{K}+H=2(\mathfrak{K}+H')$, und somit ist $\mathfrak{K}+H'$ ambig; gilt ebenso $\mathfrak{K}'=\mathfrak{K}+H'$, so ist $\mathfrak{K}+H$ ambig. Daher muss $\mathfrak{L}$ unter den Klassen $L$, $L'$, $L''$ usw. vorkommen. Unter den drei Klassen $\mathfrak{K}$, $\mathfrak{K}+H$, $\mathfrak{K}+H'$ kann aber offenbar nicht mehr als eine ambig sein. Wären nämlich sowohl $\mathfrak{K}$ als auch $\mathfrak{K}+H$ ambig, also ihren jeweiligen entgegengesetzten Klassen $\mathfrak{K}'$, $\mathfrak{K}'+H'$ gleich, so wäre $\mathfrak{K}+H=\mathfrak{K}+H'$. Dasselbe folgt aus der Annahme, $\mathfrak{K}$ und $\mathfrak{K}+H'$ seien ambig; wären schließlich $\mathfrak{K}+H$, $\mathfrak{K}+H'$ ambig, also ihren entgegengesetzten Klassen $\mathfrak{K}'+H'$, $\mathfrak{K}'+H$ gleich, so wäre $\mathfrak{K}+H+\mathfrak{K}'+H=\mathfrak{K}'+H'+\mathfrak{K}+H'$, woraus $2H=2H'$ oder $H'=H$ folgt. Daher gibt es nur eine eigentlich primitive ambige Klasse, die mit $L$ zusammengesetzt $\mathfrak{L}$ ergibt; somit sind alle $L$, $L'$, $L''$ usw. verschieden.

Die Anzahl der ambigen Klassen in einer \textit{abgeleiteten} Ordnung ist offenbar gleich der Anzahl der ambigen Klassen in der primitiven Ordnung, von der sie abgeleitet ist; sie lässt sich also nach dem Vorigen stets bestimmen.

\Needspace{8\baselineskip}
\Article{260}
\textsc{Aufgabe.} \textit{Eine eigentlich primitive Klasse $K$ der Determinante $D$ entsteht durch Verdopplung einer eigentlich primitiven Klasse $k$ derselben Determinante. Gesucht sind a\,l\,l\,e solchen Klassen, deren Verdopplung die Klasse $K$ ergibt.}

\textit{Lösung.} Es sei $H$ die Hauptklasse der Determinante $D$, und $H'$, $H''$, $H'''$ usw. seien die übrigen eigentlich primitiven ambigen Klassen derselben Determinante. Die Klassen, die aus ihrer Komposition mit $k$ entstehen, nämlich $k+H'$, $k+H''$, $k+H'''$, seien mit $k'$, $k''$, $k'''$ usw. bezeichnet. Dann sind alle Klassen $k$, $k'$, $k''$ usw. eigentlich primitiv von der Determinante $D$ und voneinander verschieden; ebenso leicht sieht man, dass aus der Verdopplung jeder die Klasse $K$ entsteht. Bezeichnet nun $\mathfrak{K}$ irgendeine eigentlich primitive Klasse der Determinante $D$, deren Verdopplung die Klasse $K$ ergibt, so muss sie unter den Klassen $k$, $k'$, $k''$ usw. enthalten sein. Setzen wir nämlich $\mathfrak{K}=k+\mathfrak{H}$, wobei $\mathfrak{H}$ eine eigentlich primitive Klasse der Determinante $D$ ist (Artikel 249). so gilt
\SourceLeaf{296}
$2k+2\mathfrak{H}=2\mathfrak{K}=K=2k$; daraus folgt leicht, dass $2\mathfrak{H}$ mit der Hauptklasse zusammenfällt, $\mathfrak{H}$ ambig ist, also unter $H$, $H'$, $H''$ usw. vorkommt, und $\mathfrak{K}$ unter $k$, $k'$, $k''$ usw. Diese Klassen geben daher die vollständige Lösung der Aufgabe.

Ist $D$ negativ, so ist offenbar von den Klassen $k$, $k'$, $k''$ usw. die Hälfte positiv, die andere negativ.

Jede eigentlich primitive Klasse der Determinante $D$, die durch Verdopplung irgendeiner solchen Klasse entstehen kann, entsteht also durch Verdopplung genau so vieler solcher Klassen, wie es eigentlich primitive ambige Klassen der Determinante $D$ gibt. Ist daher die Anzahl aller eigentlich primitiven Klassen der Determinante $D$ gleich $r$ und die Anzahl aller eigentlich primitiven ambigen Klassen dieser Determinante gleich $n$, so ist die Anzahl aller eigentlich primitiven Klassen derselben Determinante, die durch Verdopplung einer solchen Klasse entstehen können, gleich $\frac{r}{n}$. Dieselbe Formel ergibt sich bei negativer Determinante, wenn die Zeichen $r$, $n$ die Anzahlen \textit{positiver} Klassen bezeichnen, das erste \textit{aller} eigentlich primitiven, das zweite nur der ambigen. So ist \textit{zum Beispiel}\ für $D=-161$ die Anzahl aller positiven eigentlich primitiven Klassen gleich $16$, die Anzahl der ambigen gleich $4$; daher muss die Anzahl aller Klassen, die durch Verdopplung irgendeiner Klasse entstehen können, gleich $4$ sein. Tatsächlich haben alle im Hauptgeschlecht enthaltenen Klassen diese Eigenschaft: Die Hauptklasse $(1,0,161)$ entsteht durch Verdopplung der vier ambigen Klassen; $(2,1,81)$ durch Verdopplung der Klassen $(9,1,18)$, $(9,-1,18)$, $(11,2,15)$, $(11,-2,15)$; $(9,1,18)$ durch Verdopplung der Klassen $(3,1,54)$, $(6,1,27)$, $(5,-2,33)$, $(10,3,17)$; schließlich $(9,-1,18)$ durch Verdopplung der Klassen $(3,-1,54)$, $(6,-1,27)$, $(5,2,33)$, $(10,-3,17)$.
